%========================================================================================
\section{API REST}
\label{sec:rest}
%========================================================================================

\subsection{Points d'accès}

Nous avons construit l'API REST avec Django REST Framework~\cite{drf}. Les tableaux~\ref{tab:rest-auth} à~\ref{tab:rest-itinerary} listent ses points d'accès. Ils demandent tous un jeton d'accès (\code{Authorization: Bearer <access>}), sauf l'inscription, la connexion, le rafraîchissement et la documentation.

\begin{table}[H]
  \centering
  \small
  \caption{Authentification (\texttt{apps/accounts/urls.py}).}
  \label{tab:rest-auth}
  \begin{tabularx}{\linewidth}{@{}>{\raggedright\arraybackslash}p{2.1cm} l X@{}}
    \toprule
    \textbf{Méthode} & \textbf{Chemin} & \textbf{Rôle} \\
    \midrule
    POST & \code{/api/auth/register/} & Inscription, renvoie l'utilisateur et une paire de jetons \\
    POST & \code{/api/auth/token/} & Connexion, donne un jeton d'accès (15 min) et un jeton de rafraîchissement (1 jour) \\
    POST & \code{/api/auth/token/refresh/} & Nouveau jeton d'accès, l'ancien jeton de rafraîchissement est révoqué \\
    POST & \code{/api/auth/logout/} & Révoque son propre jeton de rafraîchissement \\
    GET & \code{/api/auth/me/} & Profil de l'utilisateur connecté \\
    \bottomrule
  \end{tabularx}
\end{table}

\begin{table}[H]
  \centering
  \small
  \caption{Voyages et participants (\texttt{apps/trips/urls.py}).}
  \label{tab:rest-trips}
  \begin{tabularx}{\linewidth}{@{}>{\raggedright\arraybackslash}p{2.1cm} l X@{}}
    \toprule
    \textbf{Méthode} & \textbf{Chemin} & \textbf{Rôle} \\
    \midrule
    GET, POST & \code{/api/trips/} & Mes voyages (20 par page) et création \\
    GET, PUT, PATCH, DELETE & \code{/api/trips/{id}/} & Détail avec destinations et activités, modification, suppression \\
    GET & \code{/api/trips/{id}/stats/} & Coût total, budget restant, coûts par catégorie \\
    POST & \code{/api/trips/{id}/clone/} & Copie du voyage en brouillon, dans une transaction \\
    GET, POST & \code{/api/trips/{id}/participants/} & Liste et ajout par courriel \\
    PATCH, DELETE & \code{/api/trips/{id}/participants/{user_id}/} & Changement de rôle et retrait \\
    \bottomrule
  \end{tabularx}
\end{table}

\begin{table}[H]
  \centering
  \small
  \caption{Destinations et activités (\texttt{apps/itinerary/urls.py}).}
  \label{tab:rest-itinerary}
  \begin{tabularx}{\linewidth}{@{}>{\raggedright\arraybackslash}p{2.1cm} l X@{}}
    \toprule
    \textbf{Méthode} & \textbf{Chemin} & \textbf{Rôle} \\
    \midrule
    GET, POST & \code{/api/destinations/} & Destinations de mes voyages et création \\
    GET, PUT, PATCH, DELETE & \code{/api/destinations/{id}/} & Détail, modification, suppression \\
    GET & \code{/api/destinations/{id}/weather/} & Météo du séjour (Open-Meteo) \\
    GET & \code{/api/destinations/{id}/country/} & Fiche du pays (Countries) \\
    GET & \code{/api/destinations/{id}/overview/} & Vue d'ensemble avec dégradation partielle \\
    GET, POST & \code{/api/activities/} & Activités de mes voyages et création \\
    GET, PUT, PATCH, DELETE & \code{/api/activities/{id}/} & Détail, modification, suppression \\
    \bottomrule
  \end{tabularx}
\end{table}

Il y a aussi \code{/api/schema/} (schéma OpenAPI), \code{/api/docs/} (Swagger UI), \code{/api/redoc/} (ReDoc), \code{/graphql/} et \code{/health/}.

\subsection{Serializers}

Les serializers transforment les modèles en JSON et valident ce que le client envoie. Ils n'acceptent que les champs d'une liste blanche. Les champs calculés ou sensibles, comme \code{owner}, \code{id}, les horodatages et les totaux, sont en lecture seule. Un client ne peut donc pas se donner un voyage en envoyant \code{"owner": 7}.

Pour le voyage, nous avons trois serializers. \code{TripSerializer} sert pour la liste, la création et la modification. \code{TripDetailSerializer} ajoute les destinations et leurs activités, imbriquées en lecture seule. Enfin, des serializers sans modèle décrivent les réponses calculées, comme \code{TripStatsSerializer}. Le mixin \code{ModelCleanMixin}, expliqué à la section~\ref{sec:validation}, applique en plus \code{Model.clean()} et les contraintes du modèle.

\subsection{ViewSets et routage}

Chaque ressource est un \code{ModelViewSet} enregistré dans un \code{DefaultRouter}, qui crée les routes de liste et de détail. Les opérations qui ne sont pas du CRUD sont des actions déclarées avec \code{@action} : \code{stats}, \code{clone}, \code{participants}, \code{weather}, \code{country} et \code{overview}.

\begin{lstlisting}[language=Python,caption={Queryset et permissions du ViewSet des voyages (\texttt{apps/trips/views.py}, extrait).},label={lst:viewset}]
class TripViewSet(TripRolePermissionMixin, viewsets.ModelViewSet):
    filterset_class = TripFilter
    search_fields = ["title", "description"]
    ordering_fields = ["start_date", "title", "created_at"]
    object_permission_by_action = {
        "update": IsTripEditor,
        "partial_update": IsTripEditor,
        "destroy": IsTripOwner,
        "participant_detail": IsTripOwner,
    }

    def get_queryset(self):
        queryset = Trip.objects.visible_to(self.request.user).with_totals().select_related("owner")
        if self.action == "retrieve":
            queryset = queryset.prefetch_related(
                Prefetch("destinations",
                         queryset=Destination.objects.order_by("order").prefetch_related("activities"))
            )
        return queryset
\end{lstlisting}

Le dictionnaire \code{object_permission_by_action} associe chaque action à sa permission objet. Si une action n'y est pas (en pratique, une lecture), il suffit d'être membre du voyage. On voit donc toutes les règles d'accès au même endroit.

\subsection{Filtres, recherche, tri et pagination}

Trois mécanismes de filtrage sont actifs par défaut. Il y a \code{django-filter}, la recherche texte de DRF et un tri stable qui départage toujours par identifiant. Les listes sont paginées par 20.

\begin{table}[H]
  \centering
  \small
  \caption{Paramètres des listes.}
  \label{tab:filtres}
  \begin{tabularx}{\linewidth}{@{}l X l@{}}
    \toprule
    \textbf{Ressource} & \textbf{Filtres} & \textbf{Tri (\texttt{ordering})} \\
    \midrule
    Voyages & \code{status}, \code{start_date__gte}, \code{end_date__lte}, \code{title__icontains}, \code{search} & date, titre, création \\
    Destinations & \code{trip}, \code{country_code}, \code{arrival_date__gte}, \code{departure_date__lte}, \code{search} & ordre, arrivée, ville \\
    Activités & \code{trip}, \code{destination}, \code{category}, \code{date}, \code{date__gte}, \code{date__lte}, \code{search} & coût, date, nom \\
    \bottomrule
  \end{tabularx}
\end{table}

Les filtres \code{?trip=} et \code{?destination=} sont de simples filtres sur un nombre, appliqués à un queryset déjà limité aux voyages de l'utilisateur. Avec l'identifiant du voyage de quelqu'un d'autre, on obtient une liste vide et pas une erreur. Le filtre ne permet donc pas de savoir si ce voyage existe.

\subsection{Codes HTTP et format des erreurs}

L'API utilise les codes HTTP pour ce qu'ils veulent dire (tableau~\ref{tab:codes}).

\begin{table}[H]
  \centering
  \small
  \caption{Codes HTTP renvoyés.}
  \label{tab:codes}
  \begin{tabularx}{\linewidth}{@{}l l X@{}}
    \toprule
    \textbf{Code} & \textbf{Code d'erreur} & \textbf{Situation} \\
    \midrule
    200, 201, 204 & (aucun) & Lecture ou modification, création, suppression \\
    400 & \code{validation_error}, \code{integrity_error} & Données invalides, contrainte violée \\
    401 & \code{not_authenticated}, \code{token_not_valid} & Jeton absent, invalide ou expiré \\
    403 & \code{permission_denied} & Membre du voyage qui n'a pas le bon rôle \\
    404 & \code{not_found} & Ressource absente ou qui appartient à quelqu'un d'autre \\
    405, 415 & \code{method_not_allowed}, \code{unsupported_media_type} & Méthode ou type de contenu refusé \\
    429 & \code{throttled} & Limite de requêtes atteinte \\
    502, 503, 504 & \code{external_service_*} & Service externe illisible, indisponible ou trop lent \\
    \bottomrule
  \end{tabularx}
\end{table}

Toutes les erreurs ont la même enveloppe. Le champ \code{code} est un identifiant stable que le programme client peut tester. Le champ \code{message} est destiné à une personne, et \code{details} donne les erreurs champ par champ.

\begin{lstlisting}[language=json,caption={Enveloppe d'erreur d'une validation.},label={lst:erreur}]
{
  "error": {
    "code": "validation_error",
    "message": "Les données envoyées sont invalides.",
    "details": {"end_date": ["La date de fin doit être postérieure ou égale à la date de début."]}
  }
}
\end{lstlisting}

C'est le gestionnaire \code{apps.core.exceptions.exception_handler} qui produit cette forme. Il convertit aussi deux exceptions qui, sans lui, donneraient une erreur 500 : la \code{ValidationError} de Django levée par les services et l'\code{IntegrityError} levée par PostgreSQL. Les URL inconnues et les erreurs inattendues en dehors de DRF (\code{handler404}, \code{handler500}) renvoient la même enveloppe, sans page HTML ni trace de pile.

\subsection{Documentation OpenAPI}

drf-spectacular~\cite{spectacular} génère un schéma OpenAPI~3 à partir des ViewSets et des serializers. Nous avons annoté chaque opération avec \code{@extend_schema} pour ajouter un résumé en français, les réponses d'erreur possibles et des exemples. L'option \code{COMPONENT_SPLIT_REQUEST} sépare les schémas de lecture et d'écriture, pour que les champs calculés n'apparaissent pas dans les corps de requête. Dans Swagger UI (figure~\ref{fig:swagger}), on peut tester l'API après avoir collé un jeton avec le bouton \emph{Authorize}.

\capture{swagger.png}{Swagger UI (\texttt{/api/docs/}). En haut, la description de l'API et le bouton \emph{Authorize}. En dessous, les opérations sur les activités.}{fig:swagger}
